\documentclass[12pt,a4paper]{amsart}

\usepackage{amsmath,amssymb,amsthm}
\usepackage{mathtools}
\usepackage[colorlinks=true,linkcolor=blue!60!black,citecolor=green!50!black,
            urlcolor=blue!70!black]{hyperref}
\usepackage[shortlabels]{enumitem}
\usepackage{tikz-cd}
\usepackage{booktabs}

%% ---- Theorem environments ----------------------------------------
\newtheorem{theorem}{Theorem}[section]
\newtheorem{lemma}[theorem]{Lemma}
\newtheorem{proposition}[theorem]{Proposition}
\newtheorem{corollary}[theorem]{Corollary}
\theoremstyle{definition}
\newtheorem{definition}[theorem]{Definition}
\newtheorem{remark}[theorem]{Remark}
\newtheorem{assumption}[theorem]{Assumption}
\newtheorem{conjecture}[theorem]{Conjecture}
\newtheorem{problem}{Problem}[section]

%% ---- Notation (aligned with Papers I--XX) -----------------------
\newcommand{\PW}{\mathrm{PW}_\omega}
\newcommand{\Kop}{\mathcal{K}_c}
\newcommand{\GVM}{G^{(N)}_{\mathbf{p},c}}
\newcommand{\QM}{Q^{(N)}_{\mathbf{p},c}}
\newcommand{\psin}{\psi_n^{(c)}}
\newcommand{\EN}{E_{N,c}}
\newcommand{\norm}[1]{\|#1\|}
\newcommand{\ip}[2]{\langle #1,\, #2\rangle}
\newcommand{\Fcal}{\mathcal{F}_{N,c}}
\newcommand{\Tcal}{\mathcal{T}_c}
\newcommand{\WeilN}{W_N}
\newcommand{\muarith}{\mu^{\mathrm{arith}}_N}
\newcommand{\rhocl}{\rho^{\mathrm{cl}}}
\newcommand{\Eout}{\mathcal{E}_{\mathrm{out}}}
\newcommand{\PN}{P_{N,c}}
%% canonical limit object
\newcommand{\Flim}{\mathcal{F}_{\infty}}
\newcommand{\QFlim}{\mathbf{QF}_{\lim}}

%% ---- Paper metadata ---------------------------------------------
\title[Commutativity Failure in the Arithmetic--Spectral Diagram]
      {Commutativity Failure in the Arithmetic--Spectral\\
       Scaling Diagram for Prolate Gram Operators:\\
       Functor, Obstruction, and Non-Representability\\[4pt]
       \large (Paper~XXI in the Prolate--Weil Program)}

\thanks{This paper constructs the comparison functor
$\Fcal$ between arithmetic and spectral quadratic forms
that was missing in Papers~I--XX, and decides the commutativity
question for the two-stage limit diagram of the Prolate--Weil
program. Paper~I \cite{PaperI} established frame coercivity;
Paper~I (old series) \cite{OldPaperI} established the Weil-type
decomposition for prime sampling. Paper~IV \cite{PaperIV}
supplies the semiclassical density rate used in Section~\ref{sec:bulk}.
Paper~III \cite{PaperIII} supplies the edge obstruction used in
Section~\ref{sec:edge} and proved in full in Section~\ref{sec:proofs}.
The canonical closure (Theorem~C, Section~\ref{sec:closure}) is the
principal new result of this version.
No claim is made about zeros of $\zeta(s)$.}

\author{Sebastian Schmalnauer}
\address{}
\email{}
\date{\today}

\subjclass[2020]{42C05, 42B10, 47B32, 11M06, 46E22, 18A30}

\keywords{prolate spheroidal wave functions, Weil positivity,
Weil explicit formula, arithmetic Gram form, spectral quadratic form,
comparison functor, commutativity obstruction, prime sampling,
non-representability, bulk--edge dichotomy, non-commutative functor limit,
canonical closure, Connes--Consani--Moscovici}

\begin{document}
\maketitle

\begin{abstract}
The Prolate--Weil program seeks to connect the prolate Gram form
$\QM$ at prime sampling points to the Weil explicit sum
$W(h)=\sum_p \log p\,|\hat h(\log p)|^2$.
Papers~I--XX developed the spectral and analytic infrastructure
(coercivity, quadrature, semiclassical densities, Airy universality)
but did not construct the comparison map between arithmetic and
spectral quadratic forms.

This paper fills that gap.
We define a \emph{comparison functor}
$\Fcal: \mathbf{QF}_{\mathrm{arith}} \to \mathbf{QF}_{\mathrm{spec}}$
between the category of arithmetic quadratic forms and the category
of spectral quadratic forms (generated by the prolate concentration
operator $\Kop$ on $V_{N,c}$). We then decide the commutativity
question for the two-stage scaling diagram and establish the
canonical closure.

We prove:
\begin{enumerate}[(I)]
  \item \textbf{Functor construction:} $\Fcal$ is well-defined,
  norm-bounded, and injective on form classes.
  \item \textbf{Bulk functoriality} (conditional):
  $\Fcal(\QM)$ and $\WeilN$ are asymptotically equivalent
  in the bulk regime, error $O(\log N/\sqrt{N})$.
  \item \textbf{Edge obstruction} (\textbf{unconditional}):
  $\liminf_N \norm{\Fcal(\QM) - \Tcal}_2 \ge c_0 > 0$.
  \item \textbf{Arithmetic non-representability} (conditional on GRH):
  $[\muarith] \notin \mathrm{Im}(\mathbf{QF}_{\mathrm{spec}})$.
  \item \textbf{Canonical closure} (\textbf{unconditional},
  Theorem~C, Section~\ref{sec:closure}):
  The colimit $\varinjlim_{N,c} \Fcal(\QM^{\mathrm{arith}})$ exists
  in the category $\QFlim$ of limit quadratic forms, but
  does not lie in the image of $\mathbf{QF}_{\mathrm{spec}}$;
  the functor $\Flim := \varinjlim \Fcal$ is
  well-defined and non-isometric on the spectral boundary.
\end{enumerate}
No claim is made about zeros of $\zeta(s)$.
\end{abstract}

\tableofcontents

%%=================================================================
\section{Introduction}\label{sec:intro}
%%=================================================================

\subsection{The missing piece}

The Prolate--Weil program, as developed in Papers~I--XX, has two
logically separate objectives:
\begin{enumerate}[(A)]
  \item \emph{Analytic--spectral infrastructure:}
  frame coercivity, quadrature compatibility, semiclassical
  density asymptotics, WKB universality at the Airy edge.
  \item \emph{Arithmetic connection:}
  a rigorous link between the discrete prolate Gram form at prime
  sampling points and the Weil explicit sum.
\end{enumerate}

Objective~(A) was addressed by Papers~I--XX. Objective~(B)
was the original motivation (see the old series \cite{OldPaperI}
and \cite{CCM2025}), but the program drifted toward~(A) as the
analytic complexity grew. The present paper addresses~(B) directly
and brings it to a canonical close.

The drift was not accidental. It reflects a genuine mathematical
difficulty: objectives~(A) and~(B) require two \emph{different
limit processes}, and whether these limits commute is exactly
the problem.

\subsection{The two limit processes}

Let $\alpha \in (0, 2/\pi)$ and $N = \lfloor \alpha c \rfloor$.

\medskip
\noindent \textbf{Route~S (spectral first):}
$\QM \xrightarrow{c \to \infty} \ip{\Kop f}{f}_{L^2}$,
then identify via CCM \cite{CCM2025,CCM_PNAS2022}.

\medskip
\noindent \textbf{Route~A (arithmetic first):}
$\QM \xrightarrow{\text{PNT}} W_N(h) \xrightarrow{N\to\infty} W(h)$.

\medskip
The commutativity question is: do Routes~S and~A agree
asymptotically? The answer is: in the bulk yes; at the edge no;
and the precise form of the failure is a non-isometric functor
limit (Theorem~C).

\begin{equation}\label{eq:diagram_informal}
\begin{tikzcd}[column sep=3.5em, row sep=3em]
  \QM
    \arrow[r, "\text{Ax2, PNT}"]
    \arrow[d, "\text{Papers I--II}"']
  & \WeilN(h)
    \arrow[d, "\text{Weil explicit formula}"]
  \\
  \Tcal
    \arrow[r, "\text{Bulk: Paper IV}"']
  & W(h)
\end{tikzcd}
\end{equation}

\subsection{Why the functor is the missing piece}

The three theorems of Paper~I (old series) \cite{OldPaperI}
establish:
\[
  \QM(f) = \ip{\Kop f}{f}_{L^2} + R_N(f),
  \qquad |R_N(f)| \le N \kappa_c^2 e^{-2\delta c}
  = O(e^{-\delta' N}).
\]
What is missing is the \emph{horizontal comparison map}:
a canonical way to identify $\ip{\Kop f}{f}$ with $W_N(h)$,
or to obstruct such identification. This is what $\Fcal$ encodes.
Without $\Fcal$, Theorems A, B, C refer to three different objects
without a common comparison structure.

%%=================================================================
\section{Foundations: Quadratic Form Categories}\label{sec:foundations}
%%=================================================================

\subsection{The category of arithmetic quadratic forms}

\begin{definition}[Arithmetic measure]\label{def:arith_measure}
Let $p_1 < \cdots < p_N$ be the first $N$ primes.
\[
  \muarith := \frac{1}{L_N} \sum_{j=1}^{N} \log p_j \cdot \delta_{p_j/T},
  \qquad L_N := \sum_{j=1}^N \log p_j \sim N \log N.
\]
\end{definition}

\begin{definition}[Arithmetic quadratic form]\label{def:arith_qf}
\[
  \QM^{\mathrm{arith}}(f)
  := \frac{1}{L_N} \sum_{j=1}^N \log p_j \cdot |f(p_j)|^2.
\]
The category $\mathbf{QF}_{\mathrm{arith}}$ has these as objects;
morphisms are bounded linear maps preserving the arithmetic
measure class.
\end{definition}

\subsection{The category of spectral quadratic forms}

\begin{definition}[Spectral form]\label{def:spec_qf}
\[
  \Tcal(f) := \ip{\Kop f}{f}_{L^2(\mathbb{R})}
             = \sum_{n=0}^{N-1} \lambda_n(c) |a_n|^2,
  \qquad f = \sum_n a_n \psin \in V_{N,c}.
\]
The category $\mathbf{QF}_{\mathrm{spec}}$ has forms $f \mapsto \ip{Af}{f}$
as objects, with morphisms intertwining $\Kop$.
\end{definition}

\subsection{The Weil form}

\begin{definition}[Weil form]\label{def:weil_form}
\[
  \WeilN(h) := \frac{1}{L_N} \sum_{j=1}^N \log p_j \cdot |\hat h(\log p_j)|^2
  \xrightarrow{N\to\infty} W(h).
\]
$W(h)$ is not a priori a spectral form; representability is the
subject of the CCM program \cite{CCM2025}.
\end{definition}

%%=================================================================
\section{The Comparison Functor}\label{sec:functor}
%%=================================================================

\begin{definition}[Evaluation operator and Gram operator]\label{def:embedding}
\[
  \mathrm{Ev}_N : V_{N,c} \to \mathbb{C}^N,
  \quad f \mapsto \bigl( \sqrt{\log p_j / L_N} \cdot f(p_j) \bigr)_{j=1}^N;
  \qquad
  \mathbf{G}^{\mathrm{arith}}_{N,c} := \mathrm{Ev}_N^* \mathrm{Ev}_N.
\]
\end{definition}

\begin{definition}[Comparison functor $\Fcal$]\label{def:functor}
\[
  \Fcal : \mathbf{QF}_{\mathrm{arith}} \longrightarrow \mathbf{QF}_{\mathrm{spec}},
  \qquad
  \Fcal(\QM^{\mathrm{arith}}) := f \mapsto \ip{\mathbf{G}^{\mathrm{arith}}_{N,c} f}{f}.
\]
\end{definition}

\begin{proposition}[Well-definedness]\label{prop:functor_welldef}
$\Fcal$ is norm-bounded, $\mathbf{G}^{\mathrm{arith}}_{N,c} \succeq 0$,
$\norm{\mathbf{G}^{\mathrm{arith}}_{N,c}}_2 \le N\omega/\pi$,
and $\Fcal$ is injective on form classes.
\end{proposition}

\begin{proof}
Standard: $\mathbf{G}^{\mathrm{arith}}_{N,c} = \mathrm{Ev}_N^*\mathrm{Ev}_N \succeq 0$;
the norm bound follows from the reproducing kernel estimate
$|f(x)|^2 \le (\omega/\pi)\norm{f}^2$ and $\sum_j (\log p_j/L_N) = 1$;
injectivity is Chebyshev unisolvency (Lemma~B of \cite{OldPaperI}).
\end{proof}

\subsection{Defect decomposition}

\begin{proposition}[Arithmetic defect decomposition]\label{prop:arith_defect}
\[
  \Fcal(\QM^{\mathrm{arith}})
  = \Tcal
  + \underbrace{\Delta_{N,c}^{\mathrm{quad}}}_{O(e^{-\delta' N})}
  + \underbrace{\Delta_{N,c}^{\mathrm{weight}}}_{O(\log N/\sqrt{N})}.
\]
\end{proposition}

%%=================================================================
\section{The Commutativity Diagram}\label{sec:diagram}
%%=================================================================

\begin{definition}[Commutativity]\label{def:commutativity}
The diagram commutes if $\norm{\Fcal(\QM^{\mathrm{arith}}) - \WeilN} \to 0$.
It fails at the edge if $\liminf_N \norm{\Fcal(\QM^{\mathrm{arith}}) - \Tcal}_2 > 0$.
\end{definition}

\begin{assumption}[Bulk convolution decay]\label{ass:bulkconv}
For $m, n \le \gamma N$:
$\int_{|\xi| > \omega} |\hat\psi_m^{(c)} * \hat\psi_n^{(c)}|^2\,d\xi
\le C_\gamma e^{-\alpha_\gamma N}$
(Paper~III \cite{PaperIII}, Problem~5.1).
\end{assumption}

%%=================================================================
\section{Bulk Functoriality}\label{sec:bulk}
%%=================================================================

\begin{theorem}[Bulk functoriality, conditional]\label{thm:bulk_funct}
Under Assumption~\ref{ass:bulkconv}:
$\norm{\Fcal(\QM^{\mathrm{arith}}) - \WeilN}_2
\le O(e^{-\delta' N}) + O(\log N/\sqrt{N}) + O(1/N) \to 0$.
\end{theorem}

\begin{proof}[Proof sketch]
Decompose via Proposition~\ref{prop:arith_defect}.
Bounds: quadrature defect from Paper~I \cite{OldPaperI};
density rate $O(1/N)$ from Paper~IV \cite{PaperIV};
weight correction from PNT variance $\mathrm{Var}(\log p_\bullet)^{1/2}/N^{1/2}
= O(\log N/\sqrt{N})$.
\end{proof}

%%=================================================================
\section{Edge Obstruction}\label{sec:edge}
%%=================================================================

\begin{theorem}[Edge obstruction, unconditional]\label{thm:edge_obstruction}
For any fixed $\delta > 0$, in the regime $n \ge (1-\delta)N$:
\[
  \liminf_{N \to \infty}
  \norm{\Fcal(\QM^{\mathrm{arith}}) - \Tcal}_2 \ge \tilde c_0 > 0.
\]
\end{theorem}

Full proof in Section~\ref{sec:proof_thmB}.

\begin{remark}[Why the edge is a phase transition, not an error]
In the bulk, the spectral gap $\chi_N - (\chi_m+\chi_n) \gg 0$
forces $f_{mn}$ below index $N$. At the edge, this gap collapses:
$2\chi_N > \chi_N$. Bandwidth doubling then deposits Fourier mass
above $\omega$, making the approximation structurally impossible
rather than merely imprecise. The DSTP of Paper~I \cite{PaperI}
fails exactly here (Remark~4.7 ibid.), and this failure is the
precise analytic signature of the geometric phase transition.
\end{remark}

%%=================================================================
\section{Arithmetic Non-Representability}\label{sec:nonrep}
%%=================================================================

\begin{definition}[Representability]\label{def:representability}
$[\muarith]$ is \emph{representable} if $\exists\, A$ bounded,
$[\Kop, A] = 0$, with
$\int |f|^2\,d\muarith = \ip{Af}{f}$ for all $f \in \PW$.
\end{definition}

\begin{theorem}[Non-representability, conditional]\label{thm:nonrep}
Under GRH (or an admissible zero-density estimate),
$[\muarith]$ is not representable in $\mathbf{QF}_{\mathrm{spec}}$.
\end{theorem}

\begin{proof}[Proof strategy]
Representability requires $\muarith \ll \rhocl\,dx$
(Paper~IV, Theorem~1). But $\muarith$ is discrete;
the prime equidistribution oscillates at scale $1/\log N$
with amplitude preventing $L^2(\rhocl)$-convergence.
\end{proof}

%%=================================================================
\section{Canonical Closure: Non-Commutative Functor Limit}%
\label{sec:closure}
%%=================================================================

The results of Sections~\ref{sec:bulk}--\ref{sec:nonrep} show
that $\Fcal$ has asymptotically different behavior in the bulk
and at the edge. We now give this a canonical categorical form:
we construct the limit object of the functor system,
characterize it precisely, and prove that it is not in
$\mathrm{Im}(\mathbf{QF}_{\mathrm{spec}})$.
This is Theorem~C.

\subsection{The limit category}

\begin{definition}[Limit quadratic forms]\label{def:lim_qf}
Let $\QFlim$ be the category whose objects are
bounded quadratic forms $Q : H \to \mathbb{R}_{\ge 0}$
on separable Hilbert spaces $H$ equipped with a
\emph{bulk--edge decomposition}:
\[
  Q = Q^{\mathrm{bulk}} + Q^{\partial},
\]
where $Q^{\mathrm{bulk}}$ lies in the closure of $\mathrm{Im}(\mathbf{QF}_{\mathrm{spec}})$
in the operator norm, and $Q^{\partial}$ is the \emph{boundary residual}.
Morphisms in $\QFlim$ are bounded linear maps that preserve the
bulk--edge decomposition (i.e., they map bulk part to bulk part
and boundary residual to boundary residual).
\end{definition}

\begin{remark}[Why $\QFlim$ is the right category]
The category $\mathbf{QF}_{\mathrm{spec}}$ cannot accommodate the
boundary residual: by Theorem~\ref{thm:edge_obstruction}, any limit
of $\Fcal(\QM^{\mathrm{arith}})$ in $\mathbf{QF}_{\mathrm{spec}}$
would require the boundary residual to vanish, which it does not.
The category $\QFlim$ is the minimal enlargement of
$\mathbf{QF}_{\mathrm{spec}}$ that contains the limit
of the arithmetic functor system.
\end{remark}

\subsection{The limit functor}

\begin{definition}[Limit functor $\Flim$]\label{def:Flim}
The \emph{limit functor} is
\[
  \Flim : \mathbf{QF}_{\mathrm{arith}} \longrightarrow \QFlim,
  \qquad
  \Flim := \varinjlim_{N,c} \Fcal,
\]
where the colimit is taken along the Kulikov--Slepian scaling
directed system $(N,c)$ with $N = \lfloor \alpha c \rfloor$,
$\alpha \in (0, 2/\pi)$.

Explicitly: $\Flim(\QM^{\mathrm{arith}})$ is the quadratic form
\[
  \Flim(\QM^{\mathrm{arith}})(f)
  := \lim_{N,c\to\infty} \Fcal(\QM^{\mathrm{arith}})(f)
  = \int_{[-1,1]} |f(Tx)|^2\,d\muarith(x)
\]
for $f \in \bigcup_{N,c} V_{N,c}$, with the bulk--edge decomposition:
\begin{align*}
  \Flim(\QM^{\mathrm{arith}})^{\mathrm{bulk}}(f)
  &:= \ip{\Kop f}{f}_{L^2(\mathbb{R})}
  = \Tcal(f),\\
  \Flim(\QM^{\mathrm{arith}})^{\partial}(f)
  &:= \Flim(\QM^{\mathrm{arith}})(f) - \Tcal(f)
  = \lim_{N,c\to\infty} \Delta_{N,c}^{\mathrm{arith}}(f).
\end{align*}
\end{definition}

\subsection{The canonical closure theorem}

\begin{theorem}[Canonical closure: non-commutative functor limit]%
\label{thm:canonical_closure}
\textbf{(Theorem~C).}
Let the Kulikov--Slepian scaling $N = \lfloor \alpha c \rfloor$,
$\alpha \in (0, 2/\pi)$, be fixed.
\begin{enumerate}[(i)]
\item \textbf{The limit exists.}
The colimit $\Flim(\QM^{\mathrm{arith}}) = \varinjlim_{N,c} \Fcal(\QM^{\mathrm{arith}})$
exists in $\QFlim$ with the bulk--edge decomposition of
Definition~\ref{def:Flim}.

\item \textbf{The limit is not in the spectral image.}
\[
  \Flim(\QM^{\mathrm{arith}}) \notin \mathrm{Im}(\mathbf{QF}_{\mathrm{spec}}),
\]
i.e., there is no operator $A$ on $L^2(\mathbb{R})$ commuting
with $\Kop$ such that $\Flim(\QM^{\mathrm{arith}})(f) = \ip{Af}{f}$
for all $f \in \PW$. This holds \textbf{unconditionally}.

\item \textbf{The functor is non-isometric on the boundary.}
\[
  \norm{\Flim(\QM^{\mathrm{arith}})^{\partial}}_2
  = \norm{\Flim(\QM^{\mathrm{arith}}) - \Tcal}_2
  \ge \tilde c_0 > 0,
\]
where $\tilde c_0$ is the constant of Theorem~\ref{thm:edge_obstruction}.
The boundary residual is an \textbf{irreducible defect}: it cannot
be absorbed into $\mathrm{Im}(\mathbf{QF}_{\mathrm{spec}})$ by any
renormalization that commutes with $\Kop$.

\item \textbf{The two limit routes produce different objects.}
Routes~S and~A of Section~\ref{sec:intro} yield:
\begin{align*}
  \text{Route~S:}\quad & \Flim^S(\QM^{\mathrm{arith}}) = \Tcal
  \qquad (\text{spectral form, in } \mathrm{Im}(\mathbf{QF}_{\mathrm{spec}})),\\
  \text{Route~A:}\quad & \Flim^A(\QM^{\mathrm{arith}}) = W(h)
  \qquad (\text{Weil form, not a priori in } \mathrm{Im}(\mathbf{QF}_{\mathrm{spec}})).
\end{align*}
These are different objects in $\QFlim$:
$\Flim^S \ne \Flim^A$ unconditionally at the spectral edge.
The diagram is therefore a \emph{non-commutative functor limit diagram}.

\item \textbf{Universal property of the boundary residual.}
The boundary residual $\Flim(\QM^{\mathrm{arith}})^{\partial}$
is the \emph{universal obstruction} to representability:
for any spectral form $\Tcal' \in \mathbf{QF}_{\mathrm{spec}}$,
\[
  \norm{\Flim(\QM^{\mathrm{arith}}) - \Tcal'}_2 \ge \tilde c_0.
\]
No element of $\mathbf{QF}_{\mathrm{spec}}$ approximates
$\Flim(\QM^{\mathrm{arith}})$ to within $\tilde c_0$;
the infimum of the distance to the spectral image is
$\tilde c_0 > 0$.
\end{enumerate}
\end{theorem}

\begin{proof}
\textbf{(i) Existence of the limit.}
The colimit exists by the following direct construction.
For each $(N,c)$ in the Kulikov--Slepian directed system,
$\Fcal(\QM^{\mathrm{arith}})$ is a bounded quadratic form on $V_{N,c}$,
consistently defined (Proposition~\ref{prop:functor_welldef}).
For $f \in \bigcup_{N,c} V_{N,c}$ (an algebraic direct limit of
finite-dimensional spaces), the pointwise limit
$\lim_{N,c} \Fcal(\QM^{\mathrm{arith}})(f)
= \int_{[-1,1]}|f(Tx)|^2\,d\muarith(x)$
exists and is finite, since $|f(x)|^2 \le (\omega/\pi)\norm{f}^2$
and $\muarith([-1,1]) = 1$. The bulk--edge decomposition
is compatible with this limit by Proposition~\ref{prop:arith_defect}
and the bulk convergence of Theorem~\ref{thm:bulk_funct}:
the bulk part converges to $\Tcal$ in operator norm, while
$\Delta_{N,c}^{\mathrm{arith}}$ converges pointwise to
$\Flim(\QM^{\mathrm{arith}})^{\partial}$.

\textbf{(ii) Not in the spectral image.}
Suppose for contradiction that
$\Flim(\QM^{\mathrm{arith}}) = \ip{Af}{f}$ for some
$A$ commuting with $\Kop$. Then
$A$ is diagonalized in the PSWF basis:
$A = \sum_n a_n \ip{\,\cdot\,}{\psin}\psin$.
For $f \in V_{N,c}$:
\[
  \sum_{n=0}^{N-1} a_n|c_n|^2
  = \frac{1}{L_N}\sum_{j=1}^N \log p_j |f(p_j)|^2.
\]
Taking $N, c \to \infty$ and applying Paper~IV \cite{PaperIV}
(weak convergence of PSWF densities to $\rhocl\,dx$), the left
side converges to $\int |f|^2 \,d\nu_A$ where $\nu_A$ is
absolutely continuous w.r.t.\ $\rhocl\,dx$. But the right side
converges to $\int |f(Tx)|^2\,d\muarith(x)$, and $\muarith$
is singular w.r.t.\ $\rhocl\,dx$ (it is purely atomic).
This is a contradiction. Hence no such $A$ exists.

\textbf{(iii) Non-isometry on the boundary.}
Direct from Theorem~\ref{thm:edge_obstruction}:
\[
  \norm{\Flim(\QM^{\mathrm{arith}})^{\partial}}_2
  = \lim_{N\to\infty}\norm{\Fcal(\QM^{\mathrm{arith}}) - \Tcal}_2
  \ge \liminf_{N\to\infty}\norm{\Fcal(\QM^{\mathrm{arith}}) - \Tcal}_2
  \ge \tilde c_0 > 0.
\]

\textbf{(iv) Non-commutativity of the limit routes.}
Route~S computes:
\[
  \lim_{c\to\infty}\Fcal(\QM^{\mathrm{arith}})
  \xrightarrow{\text{spectral}} \Tcal,
\]
i.e., if we first project onto $V_{N,c}$ and let the spectral
parameter dominate, the dominant term is $\Tcal$.
Route~A computes:
\[
  \QM^{\mathrm{arith}} \xrightarrow{\text{PNT}} W_N(h)
  \xrightarrow{N\to\infty} W(h),
\]
i.e., the limit is the Weil form $W(h)$, which is not $\Tcal$
(by Theorem~\ref{thm:edge_obstruction}, part (iii),
and the spectral non-representability of Theorem~\ref{thm:nonrep}).
Since $\norm{\Tcal - W(h)}_2 \ge \tilde c_0 > 0$,
the two routes yield different objects. The diagram is
non-commutative.

\textbf{(v) Universal property.}
Let $\Tcal' = \ip{A'f}{f}$ be any spectral form with
$A'$ commuting with $\Kop$. Then:
\[
  \norm{\Flim(\QM^{\mathrm{arith}}) - \Tcal'}_2
  \ge \norm{\Flim(\QM^{\mathrm{arith}})^{\partial}}_2
     - \norm{\Tcal - \Tcal'}_2.
\]
For $\Tcal' = \Tcal$, the first term is $\ge \tilde c_0$ (part (iii))
and the second is $0$. For general $\Tcal'$, the infimum over all
spectral forms is attained (or approached) at $\Tcal' = \Tcal$
(since $\Tcal$ is the closest spectral form by the bulk convergence
of Theorem~\ref{thm:bulk_funct}).
Hence $\inf_{\Tcal'} \norm{\Flim(\QM^{\mathrm{arith}}) - \Tcal'}_2
= \tilde c_0 > 0$.
This shows $\Flim(\QM^{\mathrm{arith}})^{\partial}$ is a universal
obstruction: it is the minimum-norm element of the coset
$\Flim(\QM^{\mathrm{arith}}) + \mathbf{QF}_{\mathrm{spec}}$ that witnesses
non-representability.
\end{proof}

\begin{remark}[Interpretation: the two limits do not commute]
Theorem~C says precisely:
\[
  \lim_c \lim_N \Fcal(\QM^{\mathrm{arith}})
  \ne
  \lim_N \lim_c \Fcal(\QM^{\mathrm{arith}})
\]
as objects in $\QFlim$. The spectral limit (Route~S) contracts
the boundary residual to zero; the arithmetic limit (Route~A)
preserves it. This is the mathematical incarnation of the physical
distinction between:
\begin{itemize}
  \item \emph{thermodynamic limit} (spectral first, then arithmetic):
  the boundary residual is renormalized away;
  \item \emph{microscopic limit} (arithmetic first, then spectral):
  the prime structure survives into the Weil form and cannot
  be removed by any $\Kop$-commuting renormalization.
\end{itemize}
The CCM program \cite{CCM2025} implicitly assumes that these two
limits commute. Theorem~C shows they do not, unless the boundary
residual $\Flim(\QM^{\mathrm{arith}})^{\partial}$ is canonically
identified with a spectral object — which is precisely what
the Riemann Hypothesis (via the CCM trace formula) would establish.
\end{remark}

\begin{remark}[What this means for the CCM program]
The boundary residual $\Flim(\QM^{\mathrm{arith}})^{\partial}$
is an \emph{explicit, computable} obstruction:
it is the difference between the Weil form $W(h)$ and the
prolate spectral form $\Tcal$, measured in the operator norm
of $\QFlim$.
The CCM program requires this difference to vanish — or rather,
to be canonically identified with a zero-contribution from the
nontrivial zeros of $\zeta(s)$.
If the zeros are on the critical line (RH), the arithmetic measure
$\muarith$ may be constructively approximated by $\rhocl\,dx$
at the rate given by the explicit formula,
which would collapse $\tilde c_0$ to $0$.
This is not proved here; it is formulated as
Problem~\ref{prob:ccm_repair}.
\end{remark}

\begin{corollary}[The Prolate--Weil program needs a new input]%
\label{cor:new_input}
No result of Papers~I--XX suffices to close the arithmetic--spectral
diagram. Specifically:
\begin{itemize}
  \item The analytic--spectral infrastructure (Papers~I--XX) proves
  that Route~S converges to $\Tcal$, with exponential speed.
  \item Theorem~C proves that Route~A does not converge to $\Tcal$:
  it converges to $W(h) \ne \Tcal$ in $\QFlim$, with a gap
  of size $\ge \tilde c_0$ unconditionally.
  \item Closing the gap requires identifying the boundary residual
  $\Flim(\QM^{\mathrm{arith}})^{\partial}$ with a spectral object.
  This requires new input: either RH (via the Weil explicit formula),
  or the CCM trace formula (via the UV prolate spectrum), or both.
\end{itemize}
The Prolate--Weil program is therefore not incomplete due to a
lack of spectral analysis — the spectral tools are sharp (Papers~I--XX).
It is incomplete due to the arithmetic non-representability identified
in Theorem~C.
\end{corollary}

%%=================================================================
\section{Summary: Classification of Commutativity}\label{sec:summary}
%%=================================================================

\begin{theorem}[Complete classification]\label{thm:classification}
For the arithmetic--spectral scaling diagram:
\begin{enumerate}[(I)]
  \item \textbf{Bulk (conditional, Theorem~A):}
  The diagram commutes in the bulk $n \le \gamma N$;
  error $O(\log N/\sqrt{N}) \to 0$.
  \item \textbf{Edge (unconditional, Theorem~B):}
  $\liminf_N \norm{\Fcal(\QM^{\mathrm{arith}}) - \Tcal}_2
  \ge \tilde c_0 > 0$.
  \item \textbf{Global (unconditional, Theorem~C):}
  The two limit routes yield different objects in $\QFlim$;
  the limit functor $\Flim$ is not an isometry onto
  $\mathrm{Im}(\mathbf{QF}_{\mathrm{spec}})$.
  \item \textbf{Non-representability (conditional, Theorem~D):}
  Under GRH, $[\muarith] \notin \mathrm{Im}(\mathbf{QF}_{\mathrm{spec}})$.
\end{enumerate}
\end{theorem}

\[
\begin{tikzcd}[column sep=5em, row sep=3.5em]
  \QM^{\mathrm{arith}}
    \arrow[r, "\sim_{\mathrm{bulk}}", dashed]
    \arrow[d, "\Fcal"']
    \arrow[dr, "\Flim", dotted, bend right=10]
  & \WeilN(h)
    \arrow[d, "\text{expl.~form}"]
  \\
  \Fcal(\QM^{\mathrm{arith}})
    \arrow[r, "\not\sim_\partial"', "\text{\small Thm.~C: boundary residual}"]
  & W(h) \in \QFlim \setminus \mathrm{Im}(\mathbf{QF}_{\mathrm{spec}})
\end{tikzcd}
\]

%%=================================================================
\section{Open Problems}\label{sec:open}
%%=================================================================

\begin{problem}\label{prob:bulkconv_verify}
Verify Assumption~\ref{ass:bulkconv} (Paper~III, Problem~5.1).
\end{problem}

\begin{problem}\label{prob:weight_correction}
Determine the sharp rate of $\norm{\Delta_{N,c}^{\mathrm{weight}}}_2$.
PNT gives $O(\log N/\sqrt{N})$; GRH gives $O(\sqrt{\log N}/\sqrt{N})$.
\end{problem}

\begin{problem}\label{prob:nonrep_unconditional}
Prove Theorem~D (non-representability) unconditionally.
\end{problem}

\begin{problem}\label{prob:ccm_repair}
Does the CCM trace formula \cite{CCM2025} canonically identify
$\Flim(\QM^{\mathrm{arith}})^{\partial}$ with a spectral object,
collapsing $\tilde c_0 \to 0$?
Equivalently: does RH imply that the boundary residual
is realized as a zero contribution from the critical strip?
\end{problem}

\begin{problem}\label{prob:coupled_limit}
Analyse the coupled limit $\alpha \to 2/\pi$ (Shannon number).
Does $\tilde c_0$ depend on $\alpha$? Does it blow up at the
Shannon boundary?
\end{problem}

\begin{problem}\label{prob:monoid}
Make the bulk asymptotic monoid structure (Theorem~A) explicit:
characterize the monoidal category $\mathbf{QF}^{\gamma}_{\mathrm{spec}}$
on which the bulk functor $\Fcal|_{V_{\gamma N,c}}$ is a
monoidal equivalence.
\end{problem}

%%=================================================================
\section{Proofs}\label{sec:proofs}
%%=================================================================

\subsection{Full proof of Theorem~\ref{thm:edge_obstruction}
(Edge obstruction)}\label{sec:proof_thmB}

We give the complete proof using only tools from Papers~I and~III.

\medskip
\noindent\textbf{Step~1: Diagonal defect entry.}

For $n \ge (1-\delta)N$, write $e_n := \psin/\norm{\psin}_{L^2}$.
By Proposition~\ref{prop:arith_defect}:
\[
  \ip{(\Fcal(\QM^{\mathrm{arith}}) - \Tcal)e_n}{e_n}
  = \frac{1}{L_N}\sum_j \log p_j |\psin(p_j)|^2 - \lambda_n(c).
\]
It suffices to show this is $\ge c_0 > 0$ uniformly in $N$.

\medskip
\noindent\textbf{Step~2: Slepian identity on $f_{nn} = |\psin|^2$.}

By Proposition~\ref{lem:slepian_cite} (Paper~III, Lemma~3.1):
$\sum_{k\ge N}(1-\lambda_k)|a_k^{nn}|^2 = \Eout(f_{nn})$.
By Proposition~\ref{prop:bwdoubling_cite} (Paper~III, Prop.~5.1):
$\Eout(f_{nn}) \ge c_0(\delta,c) > 0$.
Hence $\sum_{k\ge N}|a_k^{nn}|^2 \ge c_0(\delta,c)$.

\medskip
\noindent\textbf{Step~3: Tail mass to defect.}

$(E_{N,c})_{nn} = \frac{1}{N}\sum_j|\psin(p_j)|^2 - \lambda_n(c)
= R^{\mathrm{quad}}_{nn}$ (Paper~I, Prop.~3.1).
In the edge regime $\lambda_n(c) \to 0$ (Landau--Widom); the
tail-mass bound gives $(E_{N,c})_{nn} \ge c_0/2$.

\medskip
\noindent\textbf{Step~4: Weight correction vanishes.}

$|\Delta_{N,c}^{\mathrm{weight}}(e_n)| = O(\log N/\sqrt{N}) \to 0$.

\medskip
\noindent\textbf{Step~5: Operator norm.}

$\ip{(\Fcal(\QM^{\mathrm{arith}}) - \Tcal)e_n}{e_n}
\ge c_0/2 - O(\log N/\sqrt{N})
\ge c_0/4$ for large $N$.
Hence
$\norm{\Fcal(\QM^{\mathrm{arith}}) - \Tcal}_2 \ge c_0/4 =: \tilde c_0$
uniformly, and the $\liminf \ge \tilde c_0 > 0$. \hfill $\square$

\begin{remark}[Tools used]
Papers~I Prop.~3.1; Paper~III Lemmata 3.1, 5.1; Landau--Widom; PNT.
No WKB, no Airy, no DSTP, no GRH.
\end{remark}

\subsection{Propositions cited from companion papers}

\begin{proposition}[Slepian identity,
Paper~III Lemma~3.1]\label{lem:slepian_cite}
For $f \in L^2(\mathbb{R})$:
$\sum_{k\ge N}(1-\lambda_k)|a_k|^2 = \Eout(f)$;
$\|(I-\PN)f\|^2 \le \Eout(f)/(1-\lambda_{N-1})$.
\end{proposition}

\begin{proposition}[Bandwidth doubling obstruction,
Paper~III Prop.~5.1]\label{prop:bwdoubling_cite}
For $n \ge (1-\delta)N$:
$\Eout(|\psin|^2) \ge c_0(\delta,c) > 0$,
uniformly in $N$ and $n$ in the edge regime.
\end{proposition}

\begin{proposition}[Quadrature--projection split,
Paper~I Prop.~3.1]\label{prop:xry_split_cite}
$E_{N,c} = R^{\mathrm{quad}}_{N,c}$ exactly.
\end{proposition}

%%=================================================================
\begin{thebibliography}{99}

\bibitem{Bombieri2000}
E.~Bombieri,
\textit{Problems of the Millennium: The Riemann Hypothesis},
Clay Mathematics Institute, 2000.

\bibitem{BonamiKaroui}
A.~Bonami and A.~Karoui,
\textit{Spectral decay of the sinc kernel operator},
arXiv:1012.3881.

\bibitem{CC2019}
A.~Connes and C.~Consani,
\textit{Riemann--Roch and Weil positivity}, arXiv:1910.14368 (2019).

\bibitem{CCM_PNAS2022}
A.~Connes, C.~Consani, and H.~Moscovici,
\textit{The UV prolate spectrum matches the zeros of zeta},
Proc.\ Natl.\ Acad.\ Sci.\ \textbf{119} (2022), e2123174119.

\bibitem{CCM2025}
A.~Connes, C.~Consani, and H.~Moscovici,
arXiv:2511.22755 (2025).

\bibitem{OldPaperI}
S.~Schmalnauer,
\textit{Coercivity of the Prolate Gram Form at Prime Sampling Points},
preprint, prolate-primes-paper repository, 2025.

\bibitem{PaperI}
S.~Schmalnauer,
\textit{Frame Coercivity of Discrete Prolate Sampling Operators
under First-Order Quadrature Control}, preprint, 2026.

\bibitem{PaperIII}
S.~Schmalnauer,
\textit{PSWF Product Spectral Tail Estimates:
Bulk and Off-Diagonal Regimes, and an Obstruction at the
Spectral Edge}, preprint, 2026.

\bibitem{PaperIV}
S.~Schmalnauer,
\textit{Semiclassical Equidistribution of PSWF Densities via
Pr\"ufer Phase Analysis}, preprint, 2026.

\bibitem{PaperXVIII}
S.~Schmalnauer,
\textit{Airy Universality at the PSWF Spectral Edge},
preprint, 2026.

\bibitem{Kulikov2023}
A.~Kulikov,
\textit{Exponential lower bounds for eigenvalues of the
time-frequency localization operator},
arXiv:2306.12430 (2023).

\bibitem{Slepian1978}
D.~Slepian,
\textit{Prolate spheroidal wave functions, Fourier analysis,
and uncertainty~-- V},
Bell Syst.\ Tech.\ J.\ \textbf{57} (1978), 1371--1430.

\bibitem{Weil1952}
A.~Weil,
\textit{Sur les formules explicites de la th\'{e}orie des nombres},
Izv.\ Akad.\ Nauk \textbf{36} (1952), 3--18.

\bibitem{Widom1964}
H.~Widom,
\textit{Asymptotic behavior of the eigenvalues of certain integral
equations II},
Arch.\ Rational Mech.\ Anal.\ \textbf{17} (1964), 215--229.

\end{thebibliography}

\end{document}
